\documentclass[11pt]{article}
\usepackage[margin=1in]{geometry}
\usepackage{amsmath,amssymb,amsthm,booktabs,array,enumitem,microtype}
\usepackage[hidelinks]{hyperref}
\usepackage[T1]{fontenc}
\usepackage{lmodern}
\newtheorem{theorem}{Theorem}
\newtheorem{lemma}{Lemma}
\newtheorem{proposition}{Proposition}
\newcommand{\HGP}{\mathrm{HG}_{P}}
\title{The Proper Hat-Guessing Number of $K_7-e$}
\author{Matthew Protti\\\small Independent Researcher}
\date{4 September 2026\\\small Public research disclosure v0.1; independent adversarial review accepted}
\begin{document}
\maketitle
\begin{abstract}
We prove that the proper hat-guessing number of the complete graph on seven
vertices with one edge removed is twelve. Two distinct explicit constructions
prove the lower bound. The first uses two block-disjoint copies of the Witt
design $S(5,6,12)$; the second uses orbit maps from a directly regenerated
sharply five-transitive twelve-point permutation group and is genuinely
order-sensitive. Both constructions satisfy a common coordinate-line
criterion that reduces the residual game to Hall's theorem. We also record a
general disjoint completion-design theorem, an even-$n$ obstruction to
set-symmetric line-permutation twin rules inside this sufficient framework,
and a prime-admissibility theorem for the design parameters. Dependency-free
Python and independent C++ code regenerate the finite core.
\end{abstract}

\section{Introduction}
In the proper hat-guessing game on a graph $G$, an adversary assigns colours
properly from a palette of size $q$. Each vertex sees the colours of its
neighbours and guesses its own colour; the players win if at least one guess is
correct. Adriaensen et al. introduced the parameter $\HGP(G)$ and proved
\[
  \HGP(G)\le |V(G)|+\chi(G)-1.
\]
For $K_n-e$, their results leave only the two values $2n-3$ and $2n-2$.
We settle the case $n=7$.

\begin{theorem}\label{thm:main}
\[
\HGP(K_7-e)=12.
\]
\end{theorem}

The graph has seven vertices and chromatic number six, so the displayed upper
bound gives $\HGP(K_7-e)\le12$. We give two twelve-colour strategies.

\section{Twin completion and coordinate lines}
Let the missing edge join nonadjacent twins $a,b$, and let the other
$k=n-2$ vertices form an ordered clique. Put $q=2n-2$. Suppose the twins use
legal functions $\alpha,\beta$ of the injective clique-colour tuple.

Fix a clique vertex and freeze the other $k-1$ clique colours. If $D$ is the
remaining colour set, then $|D|=n+1$. Restricting the twin rules to this
coordinate line gives maps $f,g:D\to D$.

\begin{lemma}[Coordinate-line criterion]\label{lem:line}
Assume on every coordinate line that $f$ and $g$ are derangement
permutations, $f(r)\ne g(r)$ for every $r$, and $g\circ f$ has no fixed
point. Then the twin rules extend to a winning $q$-colour strategy on
$K_n-e$.
\end{lemma}

\begin{proof}
For distinct twin colours $x\ne y$, failure to cover every legal unseen clique
colour would force $f(y)=x$ and $g(x)=y$, making $y$ fixed by $g\circ f$.
For equal twin colour $x$, the two legal candidates $f^{-1}(x)$ and
$g^{-1}(x)$ are distinct. Hence, after removing colourings already covered by
the twins, every labelled clique view lies in at most $k$ residual colourings.
The residual incidence graph has left degree $k$ and right degree at most $k$;
therefore Hall's condition holds by edge counting. A matching assigns
consistent clique guesses to all residual colourings.
\end{proof}

No explicit giant matching is needed: the degree inequality proves existence.

\section{Two disjoint completion designs}
\begin{proposition}[Disjoint completion designs]\label{prop:design}
If two block-disjoint Steiner systems
\[
 S(n-2,n-1,2n-2)
\]
exist on the same colour set, then $\HGP(K_n-e)=2n-2$.
\end{proposition}
\begin{proof}
For each $(n-2)$-set $B$, let $f_i(B)$ be the unique point completing $B$ to
a block of the $i$th design. On a coordinate line each completion map is a
fixed-point-free involution. Block-disjointness makes the two maps pointwise
unequal. For involutions, $g(f(r))=r$ if and only if $f(r)=g(r)$, so the
composition is fixed-point-free. Apply Lemma~\ref{lem:line} and the general
upper bound.
\end{proof}

Our certificate contains two explicit 132-block systems $\mathcal W_A$ and
$\mathcal W_B$. Direct verification shows that every one of the
$\binom{12}{5}=792$ pentads is completed exactly once in each system and that
the systems share no block. Proposition~\ref{prop:design} proves the lower
bound in Theorem~\ref{thm:main}.

For direct confirmation, on all 495 frozen four-set lines both maps have cycle
type $(2,2,2,2)$, are pointwise unequal, and have fixed-point-free
composition. The composition has type $(4,4)$ on 330 lines and
$(2,2,2,2)$ on 165 lines.

\section{A distinct sharply transitive orbit construction}
On $\Omega=\{0,\ldots,11\}$ consider
\[
\begin{aligned}
 a&=(0\ 1\ 2\ 3\ 4\ 5\ 6\ 7\ 8\ 9\ 10),\\
 b&=(2\ 6\ 10\ 7)(3\ 9\ 4\ 5),\\
 c&=(0\ 11)(1\ 10)(2\ 5)(3\ 7)(4\ 8)(6\ 9).
\end{aligned}
\]
A breadth-first closure gives 95,040 permutations, and the images of
$(0,1,2,3,4)$ are all 95,040 ordered injective five-tuples. Thus sharp
five-transitivity is proved as a finite fact, without using a group
classification premise.

For $y\in\{5,\ldots,11\}$ define
\[
 F_y(g(0),\ldots,g(4))=g(y).
\]
The complete line classification is:
\begin{itemize}[nosep]
\item $F_6$ has type $(2,2,2,2)$ on every one of 59,400 ordered lines;
\item each of $F_5,F_7,F_8,F_9,F_{10},F_{11}$ has type $(4,4)$ on every line;
\item exactly the six pairs $\{F_6,F_y\}$ with $y\ne6$ have pointwise
unequal outputs and fixed-point-free composition on every line.
\end{itemize}
Choose $\alpha=F_6$ and $\beta=F_5$. Lemma~\ref{lem:line} again proves the
lower bound.

This proof is not merely the first proof in different notation. The map
$F_6$ is exactly the completion map of $\mathcal W_A$, but $F_5$ is
order-sensitive: on the 120 orderings of each pentad, it assumes six distinct
outputs.

\section{General structural consequences}
\begin{proposition}[Even-$n$ obstruction]\label{prop:even}
For even $n$, no set-symmetric twin rule can be a permutation on every
coordinate line within the sufficient framework of Lemma~\ref{lem:line}.
\end{proposition}
\begin{proof}
For a fixed output colour $y$, the preimages form a Steiner system
$S(n-3,n-2,2n-3)$. The number of blocks through a fixed $(n-4)$-set would be
$(n+1)/2$, which is not integral when $n$ is even.
\end{proof}
The statement is deliberately scoped: it does not characterize every possible
winning strategy.

\begin{proposition}[Prime admissibility]\label{prop:prime}
The standard divisibility conditions for $S(n-2,n-1,2n-2)$ hold if and only
if $n$ is prime. When such a design exists, its block count is the Catalan
number $C_{n-1}$.
\end{proposition}
\begin{proof}
Writing $j=t-s$, the required numbers are
\[
 \lambda_s=\frac{\binom{n+j}{j}}{j+1}
          =\frac1n\binom{n+j}{j+1}.
\]
For prime $n$, Lucas's theorem gives the required divisibility for
$1\le j\le n-2$. If $n$ is composite and $p$ is its least prime factor, take
$j=p-1$; then $\binom{n+p-1}{p-1}\equiv1\pmod p$, so the corresponding
$\lambda_s$ is not integral. The block count identity is
$\frac1n\binom{2n-2}{n-1}=C_{n-1}$.
\end{proof}
This is an admissibility result, not an existence theorem for larger primes.

\section{Finite census and verification}
For either compatible construction,
\[
(12)_5=95,040
\]
ordered clique colourings occur. The twins each have seven legal colours and,
because their guesses differ, cover $7+7-1=13$ of the 49 twin-colour pairs,
leaving 36. Hence
\[
\begin{array}{rcl}
\text{proper colourings}&=&4,656,960,\\
\text{twin-covered}&=&1,235,520,\\
\text{residual colourings}&=&3,421,440,\\
\text{residual incidence edges}&=&17,107,200.
\end{array}
\]
The maximum residual right degree is five.

The release contains a standard-library Python constructor, an independently
written C++20 constructor, a data-only verifier of the frozen tables, and
rejecting negative controls. An independent adversarial review rebuilt both
finite constructions rather than trusting the stored result summaries and
returned \texttt{ACCEPT\_K7E\_THEOREM}.

\section*{Scope and AI-use disclosure}
The result does not solve the general $K_n-e$ problem or $K_8-e$, and it does
not assert the existence of the completion designs for every prime $n$.
Conventional peer review is not claimed.

OpenAI GPT-5.6 Pro was used substantively in literature search, target
selection, construction search, proof development, code generation,
verification, adversarial critique, and manuscript preparation. Matthew
Protti selected and framed the target, directed and evaluated the work,
commissioned independent review, set the scope, approved disclosure, and
accepts responsibility.

\begin{thebibliography}{9}
\bibitem{properhat}
S. Adriaensen et al., \emph{Hat guessing with proper colorings},
arXiv:2603.04909v4, 2026.
\bibitem{kramermesner}
E. S. Kramer and D. M. Mesner, \emph{Intersections among Steiner systems},
J. Combin. Theory Ser. A 16 (1974), 273--285.
\end{thebibliography}
\end{document}
